\section{Worked problems: re-reads, records and a long prefill}

\subsection{Where the re-reads go}
\textbf{Problem.} One head of causal attention at $S = 8{,}192$ and $d_h = 128$ in
BF16 runs on an H100 (989~TFLOP/s dense BF16, 3.35~TB/s of HBM) with
$128 \times 128$ tiles. How long do its FLOPs take at peak? How long would its
memory traffic take if every tile loaded its keys and values from HBM, and if
the unfused kernel ran instead, computing the full square with 16-bit scores?
What must happen for the tiled kernel to reach its compute roof?

\textbf{Solution.} The kernel computes 2,080 tiles, the 64 on the diagonal in
full before masking them, and each costs $4 \times 128^3 = 8.39$~MFLOP: 17.4~GFLOP
in all, 17.6~$\mu$s at peak. If every tile loaded its 32~KiB of keys and 32~KiB
of values from HBM, the traffic would be 136.3~MB, plus 4.2~MB for $\mathbf{Q}$
and $\mathbf{O}$: 41.9~$\mu$s at 3.35~TB/s, so the kernel could run at no more
than 42\% of peak. The unfused kernel is worse: $8{,}192^2$ scores at 8 bytes
each over its four passes is 537~MB, 160~$\mu$s, against 34.7~$\mu$s for the
full square's arithmetic --- at most 22\% of peak (all derived). Tiling alone
is not enough. If each head's keys and values crossed HBM only once, the whole
kernel would read and write 8.4~MB, 2.5~$\mu$s, far below the 17.6~$\mu$s of
arithmetic. The 136~MB of re-reads must then come from on-chip reuse --- the
L2 cache, or on Hopper the shared memory of neighbouring blocks in a cluster
(\Chref{matmul-hardware}) --- at $136.3~\text{MB} / 17.6~\mu\text{s} = 7.7$~TB/s,
more than twice HBM's bandwidth (derived). A tile's intensity decides how much
reuse the kernel needs; the chip's reuse decides whether it gets it.

\subsection{How often does a row rescale?}
\textbf{Problem.} The lab's tiled kernel rescaled 4,499 times at 2,048 tokens and
28,383 times at 8,192. Explain the counts, treating the blocks of scores a row
meets as random and interchangeable. Then explain why a threshold of $\ln 256$
removed every rescaling on ordinary keys but only 87\% of them on the ramp keys.

\textbf{Solution.} A row rescales at the $j$-th block it meets ($j \ge 2$) only
if that block holds the largest score of the first $j$ --- for interchangeable
blocks, a chance of $1/j$. A row that meets $n$ blocks is therefore rescaled
$\sum_{j=2}^{n} 1/j = H_n - 1$ times on average, where $H_n$ is the $n$-th
harmonic number, about $\ln n + 0.58$. The 64 rows of the $i$-th query block
meet $i$ key blocks, so with $N = S/64$ blocks the expected total is
\[
64 \sum_{i=1}^{N} (H_i - 1) = 64\,\bigl[(N+1)H_N - 2N\bigr].
\]
For $N = 32$ that is $64 \times (33 \times 4.0585 - 64) = 4{,}476$; for $N = 128$,
$64 \times (129 \times 5.4331 - 256) = 28{,}472$ (derived) --- within 1\% of the
measured counts. Records grow only logarithmically, which is why 28,383
rescalings are 5\% of the kernel's 528,384 row-block updates. (At 256 tokens the
formula gives 155 against 125 measured: the diagonal block, which each row sees
only up to itself, holds fewer records than a full block, and at four blocks per
row it matters more.)

The threshold asks for more than a new record: a rise of more than
$\ln 256 = 5.55$. The lab's queries and keys are uniform on $[-2, 2]$, so each of
a score's 64 products has variance $(4/3)^2$, and after scaling by $1/8$ the
score's standard deviation is $4/3$ (derived). The largest of the 64 scores in a
row's first block is typically about 3; a later block would need a score near
8.7, six and a half standard deviations out, and among the 2.1~million scores
at 2,048 tokens that never happened (measured: no rescalings). The ramp keys scale key $j$ by $1 + 4j/S$, so
the spread of the scores grows fivefold along the sequence and a row's maximum
keeps climbing. A climbing maximum still crosses the threshold, but only once
per 5.55 of climb, instead of at every small new record: 1,604 rescalings
instead of 12,456 (measured). Either way the answer is exact; the threshold only
bounds how large the exponentials may grow in between.

\subsection{Predicting a long prefill}
\textbf{Problem.} The 20-bytes-per-score model fitted nothing and matched
llama.cpp's unfused prefill of 2,048 and 4,096 tokens. Before measuring,
predict the extra time of the unfused path at 8,192 tokens, and the compute
buffer it would reserve at 32,768 cells. Then compare with the measurement.

\textbf{Solution.} 8,192 tokens are 16 micro-batches of 512, so the prefill
writes $176{,}160{,}768 \times (16 \times 17)/2 = 24.0$~billion scores. At 20 bytes
each that is 479~GB, which at 49~GB/s takes 9.78~s (derived). Measured, two
repetitions each: 51.79~s with flash attention off and 41.74~s with it on, a
difference of 10.05~s --- 3\% more than predicted --- and the unfused path's
throughput fell from 225 to 158 tokens per second between 4,096 and 8,192
tokens while the fused path's fell from 261 to 197. The buffer grows by 400~MiB
per 8,192 cells from 416.01~MiB at 8,192, so at 32,768 cells it would be
$816.01 + 2 \times 400 = 1{,}616$~MiB (derived); the same model puts the unfused
path's extra time for a 32,768-token prompt at about 150~s. A model that has
been right three times, with no constant fitted to it, is a model worth planning
with: on this machine the unfused path's penalty grows with the square of the
prompt --- 0.73~s at 2,048 tokens, 10~s at 8,192, about 150~s at 32,768.
